\section{6. Integralrechnung}

\subsection{Stammfunktion, unbestimmtes Integral}

\Def $I \subset \RR$ Intervall, $f: I \to \RR$. $F: I \to \RR$ diff.bar mit Eigenschaft
$F'(x) = f(x) \forall x \in \RR$ nennt man \key{Stammfunktion} von $f$ auf $I$. Das 
Bestimmen von $F(x)$ aus $f(x)$ nennt man (\key{unbestimmte}) \key{Integration}

\Def Die Familie Funktionen $F(x) + C$ (wobei $F$ Stammfunktion vom \key{Integrand} $f$ ist)
nennt man \key{unbestimmtes Integral} und man schreibt
$$\int f(x) dx$$

\Satz \key{Rechenregeln}: $f,g$ Funktionen, $k \in \RR$. Folgendes gilt:
\vspace{-0.6em}

\begin{tabular}{llll}
    $\displaystyle \int \hspace{-2pt} f(x) \pm g(x)dx = \int \hspace{-2pt} f(x) dx \pm \int \hspace{-2pt} g(x) dx$ & \hspace{-14pt}
    $\displaystyle \int \hspace{-2pt} \frac{1}{x} dx = \ln(|x|) + C$ \\
    
    $\displaystyle \int \hspace{-2pt} k f(x) dx = k \int \hspace{-2pt} f(x) dx$ & \hspace{-14pt}
    $\displaystyle \int \hspace{-2pt} \sin(x) dx = -\cos(x) + C$ \\
    
    $\displaystyle \int \hspace{-2pt} k dx = kx + C$ & \hspace{-14pt}
    $\displaystyle \int \hspace{-2pt} \cos(x) dx = \sin(x) + C$ \\
    
    $\displaystyle \int \hspace{-2pt} x^n dx = \frac{1}{n+1}x^{n+1} + C \text{ für } n \ne -1$ & \hspace{-14pt}
    $\displaystyle \int \hspace{-2pt} \sinh(x) dx = \cosh(x) + C$ \\
    
    $\displaystyle \int \hspace{-2pt} e^x dx = e^x + C$ & \hspace{-14pt}
    $\displaystyle \int \hspace{-2pt} \cosh(x) dx = \sinh(x) + C$
\end{tabular}

\Satz \key{Partielle Integration}: $f, g$ Funktionen.
$$\int \underbrace{f(x)}_{\downarrow}\cdot\underbrace{g'(x)}_{\uparrow} dx = f(x) \cdot g(x) - \int f'(x) \cdot g(x) dx$$

\Satz \key{Integration durch Substitution}: $f, g$ Funktionen.
$$\int f'(\underbrace{g(x)}_y) \cdot \underbrace{g'(x) dx}_{dy} = \int \frac{df}{dy} dy = f(y) + C = f(g(x)) + C$$

\Satz $\star$ \key{Partialbruchzerlegung} (PBZ): für Polynome $p(x)$, $q(x)$ und $f(x) = \frac{p(x)}{q(x)}$
kann man $\displaystyle \int f(x) dx$ wie folgt Lösen:
\vspace{-0.6em}
\begin{enumerate}
    \item Falls $\deg(p) \ge \deg(q)$: Polynomdivision $\to$ $f(x) = s(x) + \frac{r(x)}{q(x)}$
    \item $q(x)$ in Linear- und (irreduzible) quadratische Faktoren zerlegen
    \item Für jeden Faktor von $q(x)$ addiert man entsprechende Terme im Ansatz $g(x)$:
    \begin{itemize}
        \item Einfache reelle Nullstelle $(x - x_1)$: $\frac{A}{x - x_1}$
        \item Mehrfache reelle Nullstelle $(x -x_1)^m$: $\frac{A_1}{x - x_1} + \cdots + \frac{A_m}{(x - x_1)^m}$
        \item quadratischer Faktor $(x^2 + ax + b)$ mit $a^2 - 4b < 0$: $\frac{Bx + C}{x^2 + ax + b}$:
    \end{itemize}
    \item Setze $\frac{r(x)}{q(x)} = g(x)$
    \item Löse für $A, B, C, \ldots$ (Nullstellen einsetzen oder Gleichungsystem)
    \item Löse $\displaystyle \int g(x) dx$
\end{enumerate}

\Satz $\star$ \key{Mehr Rechenregeln} für PBZ: $a, b \in \RR$ beliebig. Folgendes gilt
\vspace{-0.6em}
\begin{itemize}
    \item $\displaystyle \int \frac{a}{x - b} dx = a \ln(|x - b|) + C$
    \item $\displaystyle \int \frac{a}{(x - b)^2} dx = -\frac{a}{x-b} + C$
    \item $\displaystyle \int \frac{1}{(x + a)^2 + b^2} dx = \frac{1}{b} \arctan(\frac{x + a}{b}) + C$
\end{itemize}

\Satz $\star$ \key{Weierstrass-Substitution}: Nach der Substit. $t = \tan(\frac{x}{2})$ gilt:
$$\sin x = \frac{2t}{1 + t^2} \hspace{10pt} \cos x = \frac{1 - t^2}{1 + t^2} \hspace{10pt} dx = \frac{2}{1 + t^2} dt$$

\begin{minipage}{\linewidth}
\Satz Folgendes kann hilfreich sein:

\begin{tabular}{l l l}
    Integrand Form & Substitution & Identität \\
    \hline
    $\sqrt{a^2 - x^2}$ & $x = a \sin \theta$ & $1 - \sin^2 \theta = \cos^2 \theta$ \\
    $\sqrt{a^2 + x^2}$ & $x = a \tan \theta$ & $1 + \tan^2 \theta = \sec^2 \theta$ \\
    $\sqrt{x^2 - a^2}$ & $x = a \sec \theta$ & $\sec^2 \theta - 1 = \tan^2 \theta$ \\
    $\tan^2 x$         &                     & $\sec^2 \theta - 1 = \tan^2 \theta$ \\
    $\cot^2 x$         &                     & $\csc^2 \theta - 1 = \cot^2 \theta$
\end{tabular}
\end{minipage}

\subsection{Bestimmtes Integral (Riemann-Integral)}

\Def \key{Zerlegung} von $[a, b]$ ist endliche Menge von Punkten
$a = x_0 < x_1 < \cdots < x_n = b$ mit $n \in \NN$. $x_i$ sind \key{Unterteilungspunkte}. \\
Führt zu \key{Partition} des Intervalls:
$[a, b] = \{x_0\} \cup (x_0, x_1) \cup \{x_1\} \cup \cdots \cup (x_{n-1}, x_n) \cup \{x_n\}$

\Def Zerlegung $a = y_0 < \cdots < y_m = b$ ist \key{Verfeinerung} der Zerlegung
$a = x_0 < \cdots < x_n = b$ falls $\{x_0, \ldots, x_n\} \subset \{y_0, \ldots, y_m\}$

\Def $f: [a, b] \to \RR$ \key{Treppenfunktion}, falls Zerlegung $a = x_0 < \cdots < x_n = b$
von $[a, b]$ existiert, s.d. für $k = 1, \ldots, n$ die Restriktion von $f$ auf $(x_{k-1}, x_k)$
$f\Big|_{(x_k-1, x_k)}$ konstant. \\
Wir sagen auch, dass $f$ eine Treppenfunktion bezüglich dieser Zerlegung ist.

\Def $f: [a, b] \to \RR$ Treppenfunktion bezüglich $a = x_0 < \cdots < x_n = b$. das
\key{Integral der Treppenfunktion} $f$ über $[a, b]$ ist die reelle Zahl
$$\int_a^b \hspace{-4pt} f(x) dx := \sum_{k=1}^{n} c_k(x_k - x_{k-1})$$
wobei $c_k$ der Wert von $f$ auf dem Intervall $(x_{k-1}, x_k)$ ist

\Def $f: [a,b] \to \RR$, $\mathcal{TF}$ Menge der Treppenfunktionen.
die Menge der \key{Obersummen} $\mathcal U (f)$ oder der \key{Untersummen}
$\mathcal L (f)$ sind:
$$\mathcal U (f) := \left\{\int_a^b \hspace{-4pt}s(x)dx \, | \, s \in \mathcal{TF}; s \ge f\right\}$$

$$\mathcal L (f) := \left\{\int_a^b \hspace{-4pt}s(x)dx \, | \, s \in \mathcal{TF}; s \le f\right\}$$

\Def $f: [a, b] \to \RR$ beschränkt ist \key{Riemann-integrierbar},
falls $\sup \mathcal L (f) = \inf \mathcal U (f)$ \\
In diesem Fal wird der Wert das \key{Riemann-Integral von $f$} genannt:
$$\int_a^b \hspace{-4pt}f(x) dx := \sup \mathcal L (f) = \inf \mathcal U (f)$$
$a$ und $b$ sind \key{untere}/\key{obere Integrationsgrenze} und $f$ den \key{Integranden}

\Satz \key{Eigenschaften des Riemann-Integral}: $f,g: [a,b] \to \RR$ (Riemann-) integrierbar,
$\alpha, \beta \in \RR$:
\vspace{-0.6em}
\begin{itemize}
    \item $\displaystyle \int_a^b \hspace{-4pt} (\alpha f + \beta g)(x) dx = \alpha \int_a^b \hspace{-4pt} f(x) dx + \beta \int_a^b \hspace{-4pt} g(x) dx$ (Linearität)
    \item $\displaystyle f \le g \implies \int_a^b \hspace{-4pt} f(x) dx \le \int_a^b \hspace{-4pt}g(x) dx$ (Monotonie)
    \item $\displaystyle \left|\int_a^b \hspace{-4pt} f(x) dx\right| \le \int_a^b \hspace{-4pt} |f(x)| dx$ (Dreiecksungleichng)
    \item $\displaystyle \int_a^b \hspace{-4pt} f(x) dx = -\int_{b}^{a} f(x) dx$ (Umkehrung der Integrationsrichtung)
    \item $\displaystyle a \le c \le b \implies \int_a^b \hspace{-4pt} f(x) dx = \int_a^c \hspace{-4pt} f(x) dx + \int_c^b \hspace{-4pt} f(x) dx$ (Aufteilung des Integrationsbereich)
\end{itemize}

\Satz Jede monotone Funktion ist Riemann-integrierbar

\Satz Jede stetige Funktion ist Riemann-integrierbar

\Def $D \subset \RR$, $(f_n)_{n \in \NN_0}$ Folge von Funkt. $f_n: D \subset \RR \to \RR$
und $f: D \subset \RR \to \RR$ Funktion.
Dann \key{konvergiert die Folge $(f_n)_{n\in\NN_0}$ punktweise gegen $f$}, falls, für jedes $x \in D$,
$(f_n(x))_{n\in\NN_0}$ gegen $f(x)$ konvergiert. $f$ heisst \key{punktweiser Grenzwert}
von $(f_n)_{n\in\NN_0}$

\Def $D \subset \RR$, $(f_n)_{n\in\NN_0}$ Folge von Funkt. $f_n: D \to \RR$ und
$f: D \to \RR$. \\
\key{Folge konvergiert gleichmässig gegen $f$}, falls
$\forall \varepsilon > 0 \, \exists N \in \NN_0 \, \forall n \ge N \, \forall x \in D : |f_n(x) - f(x)| < \varepsilon$

\Satz $D \subset \RR$, $(f_n)_{n \in \NN_0}$ Folge stetiger Funkt. $f_n: D \to \RR$, die gleichmässig
gegen $f: D \to \RR$ konvergiert. Dann $f$ stetig

\Satz $(f_n)_{n \in \NN_0}$ Folge int.barer Funkt. $f_n: [a, b] \to \RR$, die gleichmässig
gegen $f: D \to \RR$ konvergiert. Dann $f$ int.bar und
$$\displaystyle \int_a^b \hspace{-4pt}f(x)dx = \lim_{n\to\infty}\int_a^b \hspace{-4pt}f_n(x)dx$$

\Satz \key{Mittelwertsatz der Integralrechnung}: Falls $f(x)$ auf $[a, b]$ stetig, gilt für ein
$c \in [a, b]$:
\vspace{-10pt}
$$f(c) = \frac{1}{b-a}\int_a^b \hspace{-4pt}f(x)dx$$
Wobei $f(c)$ der \key{Mittelwert} von $f$ über $[a, b]$ ist

\Satz \key{Hauptsatz der Integral- und Differentialrechnung}: $f: [a, b] \to \RR$ stetig.
Dann ist für alle $C \in \RR$ $F: [a, b] \to \RR$ Stammfunktion von $f$:
$$F(x) := \int_{a}^{x} \hspace{-4pt} f(y) dy + C$$
Ausserdem ist jede Stammfunktion von der obigen Form

\Kor $F: [a, b] \to \RR$ stetig diff.bar. Dann für alle $x \in [a, b]$
$$F(x) = F(a) + \int_{a}^{x} \hspace{-4pt}F'(t) dt$$

\Kor $f: [a, b] \to \RR$ stetig und $F: [a, b] \to \RR$ Stammfunktion von $f$.
$$\int_a^b \hspace{-4pt} f(t) dt = F(x)\Big|_a^b = F(b) - F(a)$$

\Satz \key{Partielle Integration}: $f,g: [a, b] \to \RR$ stetig diff.bar.
$$\int_a^b \hspace{-4pt} f(x)g'(x) dx = f(x)g(x) \Big|_a^b - \int_a^b \hspace{-4pt} f'(x)g(x) dx$$

\Satz \key{Substitution, Version 1}: $I, J \subset \RR$ Intervalle, $f: I \to J$ stetig diff.bar und
$g: J \to \RR$ stetig. Dann für alle $[a, b] \subset I$:
$$\int_a^b \hspace{-4pt} g(f(x))f'(x) dx = \int_{f(a)}^{f(b)} \hspace{-4pt} g(y) dy$$

\Satz \key{Substitution, Version 2}: $I, J \subset \RR$ Intervalle, $f: I \to J$ stetig diff.bar
und $g: J \to \RR$ stetig. Für $[a, b] \subset I$ gilt $f'(x) \ne 0$ für alle $x \in [a, b]$
und $f^{-1}: [f(a), f(b)] \to \RR$ die Inverse von $f\Big|_{[a,b]}$. Dann gilt
$$\int_a^b \hspace{-4pt} g(f(x)) dx = \int_{f(a)}^{f(b)} \hspace{-4pt} g(y)(f^{-1})'(y)dy$$

\Satz $\star$ Für $p > 0$ gilt:
$$\int_{0}^{1} \frac{1}{x^p} dx \text{ konvergiert} \iff p < 1$$

\Satz $\star$ Für $p > 0$ gilt:
$$\int_{1}^{\infty} \frac{1}{x^p} dx \text{ konvergiert} \iff p > 1$$
